% =============================================================================
% Ethics Considerations
% S&P requires a separate, well-marked section at the end of the paper.
% Placed before references; does not count toward the page limit.
% =============================================================================

\section*{Ethics Considerations}

\paragraph{Data and human subjects.}
This work studies the privacy risk of datasets and involves no new
collection of or interaction with human subjects. All experiments use the
31 publicly available, previously released datasets listed in
Appendix~\ref{app:datasets} (including Adult, COMPAS, Purchase100,
Texas100, NHANES, MovieLens-1M, and Gowalla).
Our use constitutes secondary analysis of already-public data. It does not
meet the definition of human-subjects research requiring institutional
review board (IRB) approval at our institutions. We collected no new
personal data and did not attempt to re-identify any individual.

\paragraph{Sensitive domains.}
Several datasets derive from sensitive contexts: criminal-justice records
(COMPAS), hospital discharge records (Texas100), clinical measurements
(Pima Diabetes), and individual mobility traces (Gowalla). We use these
datasets solely to measure membership inference risk, never to make
any inference, prediction, or decision about the individuals they describe.
We are aware that some of these datasets, most notably COMPAS, carry
documented demographic biases. Because we do not train or release deployed
predictors and draw no conclusions about individuals, our use does not
propagate those biases into downstream decisions.

\paragraph{Dual use.}
Our central artifact, DPRI, is dual-use. Because it estimates a dataset's
privacy risk before training, an adversary with access to candidate
datasets could in principle use it as a reconnaissance tool to prioritize
high-risk targets (Section~\ref{sec:implications}). We make three
observations that, on balance, support release. First, DPRI is computed
only from intrinsic statistical properties of a dataset and confers no
capability to attack any specific deployed model or service. Our
experiments attack no live system, so no vendor vulnerability disclosure is
applicable. Second, any adversary who can compute DPRI can already mount the
membership inference attacks we study, so DPRI provides at most a marginal
prioritization advantage. Third, the defensive value is substantial and
asymmetric: DPRI lets data holders anticipate privacy risk at collection
time and calibrate protections (e.g., differential privacy budgets) before
any model is released. We therefore frame and release DPRI primarily as a
defensive auditing tool.

\paragraph{Competing interests.}
The authors declare no competing financial or non-financial interests in
relation to this work.
